\scanpage{49}
\mspage{48}
\begin{center}- 13 -\end{center}
\hfill n$^\circ$;276

de la proposition précédente il suffit pour cela de faire voir que l'homomorphisme associé $G(u)$ est surjectif, autrement dit que pour tout $x\in E_n$ il existe des $a_i\in A_{n-p_i}$ tels que l'on ait $x\equiv a_ix_i\pmod{E_{n+1}}$ ; mais cela signifie précisément que les $\tilde x_i$ engendrent le $G(A)$-module $G(E)$ : en effet, si $\tilde x\in G_n(E)$, on peut trouver des $\tilde a_i\in G(A)$ tels que $\tilde x=\sum\tilde a_i\tilde x_i$, et cette relation subsiste si l'on remplace $\tilde a_i$ par sa composante de degré $n-p_i$.

\uline{Corollaire} 1 -- \uline{Soit} $E$ \uline{un module filtré séparé sur un anneau filtré complet} $A$; \uline{Si} $G(E)$ \uline{est} \inserted{un} $G(A)$-\uline{module de type fini} (resp. \uline{noethérien}) \uline{alors} $E$ \uline{est un} $A$-\uline{module de type fini} (resp. \uline{noethérien}).

Soit en effet $F$ un sous-module de $E$, et munissons le de la filtration induite ; $G(F)$ est alors un sous-module homogène du $G(A)$-module gradué $G(E)$ ; si $G(F)$ est de type fini, on peut l'engendrer à l'aide d'éléments homogènes en nombre fini -- il suffit pour le voir de prendre les composantes homogènes d'un système fini de générateurs de $G(F)$ ; en considérant des;éléments de $F$ dont ces générateurs sont les formes initiales on voit, à l'aide de la Proposition 7, que $F$ est lui-même un module de type fini, d'où immédiatement le corollaire.

\uline{Corollaire} 2 -- \uline{Soit} $(A,A_n)$ \uline{un anneau filtré séparé et complet} ; \uline{si} $G(A)$ \uline{est un anneau noethérien à gauche}, $A$ \uline{est noethérien à gauche}.

Il suffit d'appliquer le corollaire précédent à $E=A$.

\uline{Corollaire} 3 -- \uline{Soient} $A$ \uline{un anneau}, $\mathrm m$ \uline{un idéal bilatère de} $A$, \uline{et} $E$ \uline{un} $A$-\uline{module}. \uline{Supposons} $A$ \uline{complet et} $E$ \uline{séparé pour la topologie} $\mathrm m$-\uline{adique}. \uline{Si des éléments} $x_i\in E$, \uline{en nombre fini, engendrent le} $A$-\uline{module} $E/\mathrm m.E$, \uline{ils engendrent} $E$.

D'après la Proposition 7, il suffit de montrer que les éléments $\tilde x_i\in E/\mathrm m E=G_0(E)$ engendrent le $G(A)$-module $G(E)$, autrement dit que, pour
